\section{Subclasses of Toeplitz BW Form that are SoS}
\label{sec:positive-subclasses}

The failure of SoS representability at large orders does not persist
under every structural restriction.  In this section, we show that 
the Toeplitz B\"ottcher--Wenzel form is SoS at every order whenever
one factor is symmetric or skew-symmetric. 
The other factor remains an arbitrary real Toeplitz matrix.


The proof proceeds in two steps.  First, \Cref{lem:positive-split} shows
that the Toeplitz B\"ottcher--Wenzel form with one symmetric factor can
be split into a term involving two symmetric Toeplitz matrices and a
term involving a symmetric and a skew-symmetric matrix.
\Cref{prop:same-symmetry} proves that the first term is SoS, while
\Cref{lem:positive-mixed} proves the same for the second.  The
skew-symmetric case follows analogously, with the roles of the symmetric
and skew-symmetric parts interchanged. 
After completing the proof, we discuss the exterior-product
interpretation, and its relation to the two-matrix De Smet--Dillen--Verstraelen--Vrancken (DDVV) inequality. 
For notational simplicity, throughout this section we write $F(X,Y)$
for $F_N(x,y)$ in \eqref{eq:FN}, where $X$ and $Y$ are the 
matrices of size $N$.  

% For real matrices $U,V$ of the same order, write
% \[
%  F(U,V)
%  :=
%  2\norm{U}_F^2\norm{V}_F^2
%  -2\langle U,V\rangle_F^2
%  -\norm{[U,V]}_F^2.
% \]
% Thus $F_N(x,y)=F(X,Y)$ and $F(U,V)=F(V,U)$.


\paragraph{Splitting the factor}
\label{subsec:positive-split}
Clearly, every real Toeplitz matrix $X$ admits an orthogonal
decomposition into a symmetric Toeplitz matrix and a skew-symmetric
Toeplitz matrix:
\begin{equation}\label{eq:positive-split}
 X=S_x+K_x,
 \qquad
 S_x=\frac{X+X^\top}{2},
 \qquad
 K_x=\frac{X-X^\top}{2}.
\end{equation}
Here $S_x$ is symmetric Toeplitz and $K_x$ is skew-symmetric Toeplitz.
Moreover, they are orthogonal with respect to the Frobenius inner
product,
$
 \langle S_x,K_x\rangle_F=0,
$
and hence
$
\|X\|_F^2=\|S_x\|_F^2+\|K_x\|_F^2.
$
The following lemma shows that this decomposition also splits the
B\"ottcher--Wenzel form.

\begin{lemma}[The Splitting Lemma]
\label{lem:positive-split}
Let $X$ be a real Toeplitz matrix, and let $S$ and $K$ be real
symmetric and skew-symmetric matrices, respectively, of the same order.
Then
\begin{equation}\label{eq:positive-split-identities}
 \begin{aligned}
 F(X,S)=F(S,X)&=F(S,S_x)+F(S,K_x),\\
 F(X,K)=F(K,X)&=F(K,S_x)+F(K,K_x),
 \end{aligned}
\end{equation}
where $S_x$ and $K_x$ are defined in \eqref{eq:positive-split}.
\end{lemma}

\begin{proof}
Symmetric and skew-symmetric matrices are orthogonal with respect to
the Frobenius inner product.  In particular,
\[
 \langle S,K_x\rangle_F=0,
 \qquad
 \|S_x+K_x\|_F^2
 =
 \|S_x\|_F^2+\|K_x\|_F^2.
\]
Moreover, $[S,S_x]$ is skew-symmetric, whereas $[S,K_x]$ is symmetric.
Hence
\[
 \|[S,S_x+K_x]\|_F^2
 =
 \|[S,S_x]\|_F^2+\|[S,K_x]\|_F^2.
\]
Substitution into the definition of $F$ proves the first identity.
For the second identity, use
$
\langle K,S_x\rangle_F=0
$
and the fact that $[K,S_x]$ is symmetric, while $[K,K_x]$ is
skew-symmetric.  The same argument then gives the second identity.
\end{proof}

\paragraph{The two types of pairs}
\label{subsec:positive-pairs}

The splitting lemma \Cref{lem:positive-split} reduces the theorem to two
types of BW forms.  \Cref{lem:positive-mixed} shows that the
symmetric--skew-symmetric term in \eqref{eq:positive-split-identities}
is SoS.  In fact, this result holds for any symmetric--skew-symmetric
pair, even without the Toeplitz assumption.  \Cref{prop:same-symmetry}
shows that the same-symmetry term in
\eqref{eq:positive-split-identities} is also SoS.

\begin{lemma}[Symmetric--skew-symmetric pair]
\label{lem:positive-mixed}
For every real symmetric matrix
$S=(s_{ij})_{1\le i,j\le N}$ and real skew-symmetric matrix
$K=(k_{ij})_{1\le i,j\le N}$ of order $N$,
\begin{equation}\label{eq:positive-mixed}
 F(S,K)
 =
 \norm{SK+KS}_F^2
 +
 4\sum_{\ell=1}^N
 \sum_{1\le i<j<k\le N}
 (s_{i\ell}k_{jk}-s_{j\ell}k_{ik}+s_{k\ell}k_{ij})^2.
\end{equation}
\end{lemma}

\begin{proof}

Since $\langle S,K\rangle_F=0$,
$
 F(S,K)=2\norm{S}_F^2\norm{K}_F^2-\norm{SK-KS}_F^2.
$
For a vector $v\in\mathbb R^N$, the exterior-product Lagrange
identity gives
\[
 \frac12\norm{v}^2\norm{K}_F^2-\norm{Kv}^2
 =
 \sum_{i<j<k}
 (v_i k_{jk}-v_j k_{ik}+v_k k_{ij})^2.
\]
This identity follows directly by expansion.
Apply it to each column of $S$ and sum to obtain
\[
 2\norm{S}_F^2\norm{K}_F^2
 =
 4\norm{KS}_F^2
 +
 4\sum_{\ell=1}^N\sum_{i<j<k}
 (s_{i\ell}k_{jk}-s_{j\ell}k_{ik}+s_{k\ell}k_{ij})^2.
\]
Finally, $(SK)^\top=-KS$, so
$\norm{SK}_F=\norm{KS}_F$.
The parallelogram identity therefore yields
$
 4\norm{KS}_F^2-\norm{SK-KS}_F^2
 =
 \norm{SK+KS}_F^2.
$
Combining the two identities proves
\eqref{eq:positive-mixed}.
Every entry of $SK+KS$, and every expression inside the
remaining squares, is a homogeneous quadratic polynomial
in the independent entries of $S$ and $K$.
\end{proof}

\begin{proposition}[Both symmetric or skew-symmetric pair]
\label{prop:same-symmetry}
For every $N\ge2$, $F(X,Y)$ is SoS if $X,Y$ are both symmetric
Toeplitz or both skew-symmetric Toeplitz.
\end{proposition}

\begin{proof}[Proof sketch]

In either case, the matrices are determined by their nonnegative
diagonal parameters. We use the wedge coordinates introduced in
\Cref{subsec:wedge-coordinates}. Set $m=N-1$ and reverse the nonzero
diagonal labels by writing
\[
 \hat z_{ab}=z_{N-b,N-a},
 \qquad 1\le a<b\le m.
\]
This indexing replaces the diagonal-length weights
$(N-a)(N-b)$ by $ab$.
We rewrite the Toeplitz form as
\begin{equation}\label{eq:positive-reduction}
 \begin{aligned}
 F^{\rm sym}(X,Y)
 =
 4N\sum_{b=1}^{N-1}(N-b)z_{0b}^2
 +4\Psi_{N-1}(\hat z),\qquad
 F^{\rm skew}(X,Y)
 =4\Psi_{N-1}(\hat z).
 \end{aligned}
\end{equation}
Note that the full calculation is given in
\Cref{app:same-symmetry-reduction}.
Here $\Psi_m(\hat z)$ is defined by
\begin{equation}\label{eq:positive-reduced-form}
 \Psi_m(\hat z)
 :=
 2\sum_{1\le a<b\le m}ab\,\hat z_{ab}^2
 -
 \sum_{\delta=1}^{m-1}
 \sum_{\substack{k\ge0\\2k<m-\delta}}
 \left(
   \sum_{p=k+1}^{m-\delta-k}\hat z_{p,p+\delta}
 \right)^2.
\end{equation}
Intuitively, the inner sum in
\eqref{eq:positive-reduced-form} selects a window from the sequence
$
\hat z_{1,1+\delta},\ldots,\hat z_{m-\delta,m},
$
successively removing $k$ terms from each end. 
All we need to prove next is that
$\Psi_m(\hat z(x,y))$ is SoS. For example, when $m=4$ and $\delta=1$,
the two windows give
\[
 (\hat z_{12}+\hat z_{23}+\hat z_{34})^2
 \qquad\text{and}\qquad
 \hat z_{23}^2.
\]
The essential observation is that the wedges satisfy the Pl\"ucker
relations recalled in \Cref{subsec:wedge-coordinates}. These relations
allow the cross terms in the negative window squares to be re-paired
without changing the polynomial in $x,y$. For instance,
\[
 -2\hat z_{12}\hat z_{34}
 =
 -2\hat z_{13}\hat z_{24}
 +2\hat z_{14}\hat z_{23}.
\]
Applying the prescribed re-pairing rule gives three types of
squares: window sums along a fixed index total, differences of
wedges with the same gap, and individual wedges.
\Cref{prop:positive-window-certificate} gives the
complete identity and its coefficients. The window and
difference terms have nonnegative integer weights.
After allocating the diagonal terms to these squares, the
remaining coefficient of $\hat z_{p,p+\delta}^2$, denoted
$s_\delta(p)$ in \Cref{app:positive-windows}, satisfies
$
 s_\delta(p)\ge p^2>0.
$
This estimate ensures that completing the squares leaves no
negative remainder.

Consequently, $\Psi_m$ is a sum of squares of linear combinations
of the wedges, after using their Pl\"ucker relations.
Since each wedge is quadratic in $x,y$, this is an SoS
representation in the original variables.
Equation~\eqref{eq:positive-reduction} proves both assertions.
\end{proof}


Combining the splitting in \Cref{lem:positive-split} with the two
results above gives the following all-order positive result.

\begin{theorem}[One symmetric or skew-symmetric factor]
\label{thm:one-sided}
For every $N\ge2$, the Toeplitz B\"ottcher--Wenzel form $F(X,Y)$ is SoS
whenever either $X$ or $Y$ is symmetric or skew-symmetric Toeplitz.
The other factor may be an arbitrary real Toeplitz matrix.
\end{theorem}

\begin{proof}
Suppose first that $X$ is symmetric Toeplitz.  By
\Cref{lem:positive-split},
$
 F(X,Y)=F(X,S_y)+F(X,K_y),
$
where the two terms are SoS by \Cref{prop:same-symmetry} and
\Cref{lem:positive-mixed}, respectively.  The skew-symmetric case is
identical with the two results interchanged.  Finally,
$F(X,Y)=F(Y,X)$ gives the same conclusion when the structural
assumption is imposed on $Y$.
\end{proof}

\begin{remark}
The splitting above relies on one factor having a fixed symmetry.  If
both factors are decomposed into symmetric and skew-symmetric parts,
additional cross terms appear and the four resulting BW forms no longer
give the whole expression.  Thus the argument proves the one-sided
result above, but does not extend directly to unrestricted Toeplitz
pairs.  The full expansion is recorded in
Appendix~\ref{app:positive-coupling}.
\end{remark}

\paragraph{An exterior-product regularization viewpoint}
The positive part of the B\"ottcher--Wenzel form admits a useful
regularization interpretation.  By the weighted Lagrange identity,
\[
 \|X\|_F^2\|Y\|_F^2-\langle X,Y\rangle_F^2
 =
 z^\top W_Nz,
 \qquad W_N\succ0,
\]
where $z=x\wedge y$ denotes the wedge coordinates introduced in
\Cref{subsec:wedge-coordinates}.  Since the commutator is linear in
$z$, say $[X,Y]=\mathscr C_Nz$, the generalized form
\[
 F_\tau(X,Y)
 :=
 \tau\bigl(
 \|X\|_F^2\|Y\|_F^2-\langle X,Y\rangle_F^2
 \bigr)
 -
 \|[X,Y]\|_F^2
\]
becomes
\[
 F_\tau(X,Y)
 =
 z^\top
 \bigl(\tau W_N-\mathscr C_N^*\mathscr C_N\bigr)z.
\]
Thus, for each fixed $N$, the exterior-product term acts as a
positive-definite regularizer in wedge space, and sufficiently large
$\tau$ makes the lifted quadratic form positive semidefinite.
\rev{Related interactions between regularization, global optimality, and
SoS exactness for quartic polynomial models are studied in
\cite{zhuCartis2026regularizedSOS,zhuCartis2025global,zhuCartis2024sosTaylor}.}
The B\"ottcher--Wenzel form corresponds to the fixed value $\tau=2$. 
This differs from a radial regularization such as
$\sigma\|(x,y)\|^4$: the exterior-product term vanishes whenever
$X$ and $Y$ are linearly dependent, and is therefore degenerate in
the original variables even though it is positive definite in wedge
coordinates.
\paragraph{Relation to the two-matrix DDVV form}

A closely related commutator polynomial arises from the
De Smet--Dillen--Verstraelen--Vrancken (DDVV) inequality for real
symmetric matrices \cite{lu2017,ge2024published}.  For two such
matrices, define
\[
 \mathcal D_N(X,Y)
 :=
 \bigl(\norm{X}_F^2+\norm{Y}_F^2\bigr)^2
 -2\norm{[X,Y]}_F^2.
\]
A direct expansion gives
\begin{equation}\label{eq:positive-ddvv-relation}
 \mathcal D_N(X,Y)
 =
 2F(X,Y)
 +\bigl(\norm{X}_F^2-\norm{Y}_F^2\bigr)^2
 +4\langle X,Y\rangle_F^2.
\end{equation}
Hence an SoS representation of $F$ immediately gives one for
$\mathcal D_N$.  The converse does not follow, since the last two
terms in \eqref{eq:positive-ddvv-relation} are additional squares. 
The two-matrix DDVV form itself admits an SoS representation for
arbitrary real symmetric matrices.  Thus the DDVV result does not
replace the Toeplitz argument in \Cref{prop:same-symmetry}; rather,
\eqref{eq:positive-ddvv-relation} shows precisely how the two
quartic forms differ.